% Base resume template. The pipeline reuses this preamble and macros when it
% drafts a tailored resume, so keep the structure and swap the content.
%
% Compile:  tectonic resumes/base/general.tex
%
% NOTE ON FONTS: this preamble loads `helvet` and sets \sfdefault, but on some
% TeX installations that silently does not take effect and the PDF renders in
% Latin Modern Roman instead. It looks fine either way — just be aware the
% output may not be the sans-serif you expect. If you want guaranteed bold and
% a predictable font, add \usepackage[T1]{fontenc} and \usepackage{lmodern}.

\documentclass[letterpaper,10pt]{article}

\usepackage[margin=0.5in,top=0.25in,bottom=0.25in]{geometry}
\usepackage[usenames,dvipsnames]{xcolor}
\usepackage{enumitem}
\usepackage{titlesec}
\usepackage[hidelinks]{hyperref}
\usepackage{tabularx}
\usepackage{helvet}
\renewcommand{\familydefault}{\sfdefault}

\definecolor{accent}{HTML}{1F3864}
\pagestyle{empty}
\setlength{\parindent}{0pt}

\titleformat{\section}{\large\bfseries\color{accent}}{}{0em}{}[\vspace{-6pt}\rule{\textwidth}{0.8pt}]
\titlespacing{\section}{0pt}{3pt}{1pt}

\newcommand{\entry}[4]{\vspace{2pt}%
  \textbf{#1} \hfill {\small #2}\\
  {\small\itshape #3} \hfill {\small\itshape #4}\vspace{0pt}}
\newcommand{\entrytwo}[2]{\vspace{2pt}\textbf{#1} \hfill {\small #2}\vspace{1pt}}

\setlist[itemize]{leftmargin=1.2em, itemsep=0pt, parsep=0pt, topsep=0pt, label=\textbullet}

\begin{document}

% ---------- HEADER ----------
\begin{center}
  {\LARGE\bfseries\color{accent} Your Name}\\[2pt]
  {\small City, State \,\textbar\, (000) 000-0000 \,\textbar\, \href{mailto:you@example.com}{you@example.com} \,\textbar\, \href{https://linkedin.com/in/you}{linkedin.com/in/you} \,\textbar\, \href{https://github.com/you}{github.com/you}}
\end{center}
\vspace{-8pt}

\section{Summary}
{\small Two or three sentences. Lead with years of experience and the domains
you've shipped in, then the thing that differentiates you. Tailor this per
posting --- it's the highest-leverage paragraph on the page.}

\section{Experience}
\entry{Job Title}{Mon YYYY -- Present}{Company}{Location}
\begin{itemize}\small
  \item \textbf{Lead with the outcome}, then the mechanism. Attach a number wherever one exists --- runtime cut, users served, incidents resolved, coverage gained.
  \item Name the tradeoff or constraint. That's what makes a bullet survive an interviewer asking ``why did you do it that way?''
\end{itemize}

\entry{Previous Job Title}{Mon YYYY -- Mon YYYY}{Company}{Location}
\begin{itemize}\small
  \item ...
\end{itemize}

\section{Projects \& Open Source}
\begin{itemize}\small
  \item \textbf{Project Name} (\href{https://github.com/you/project}{github.com/you/project}) --- what it does, the stack, and the hard part.
\end{itemize}

\section{Education}
\entrytwo{Degree --- Institution, Location}{Mon YYYY}\\[2pt]
\entrytwo{Earlier Degree --- Institution, Location}{Mon YYYY}

\section{Technical Skills}
{\small
\renewcommand{\arraystretch}{0.95}
\begin{tabularx}{\textwidth}{@{}l X@{}}
\textbf{Languages:} & \\
\textbf{Data \& Systems:} & \\
\textbf{Cloud \& DevOps:} & \\
\textbf{Practices:} & \\
\end{tabularx}}

\end{document}
